\documentclass[10pt,letterpaper]{article}
\usepackage[letterpaper,left=0.82in,right=0.82in,top=0.72in,bottom=0.69in,headheight=14pt,footskip=22pt]{geometry}
\usepackage{fontspec}
\IfFontExistsTF{Linux Libertine O}{\setmainfont{Linux Libertine O}}{\setmainfont{TeX Gyre Termes}}
\IfFontExistsTF{Inter}{\setsansfont{Inter}}{\setsansfont{TeX Gyre Heros}}
\usepackage[protrusion=true,expansion=false]{microtype}
\usepackage{xcolor}
\usepackage{hyperref}
\usepackage{tabularx}
\usepackage{array}
\usepackage{enumitem}
\usepackage{fancyhdr}
\usepackage{needspace}
\usepackage{ragged2e}
\usepackage{titlesec}
\usepackage{parskip}
\definecolor{ink}{HTML}{26332E}
\definecolor{accent}{HTML}{426558}
\definecolor{muted}{HTML}{66716C}
\definecolor{rule}{HTML}{D5DDD6}
\hypersetup{colorlinks=true,linkcolor=accent,urlcolor=accent,citecolor=accent,pdfauthor={Kaicheng Dong},pdftitle={Kaicheng Dong - Academic Curriculum Vitae},pdfsubject={2027 Mathematics PhD applications}}
\color{ink}
\setlength{\parindent}{0pt}
\setlength{\parskip}{4pt}
\renewcommand{\baselinestretch}{1.04}
\pagestyle{fancy}
\fancyhf{}
\renewcommand{\headrulewidth}{0pt}
\renewcommand{\footrulewidth}{0pt}
\fancyfoot[L]{\footnotesize\color{muted}\sffamily Updated November 2026}
\fancyfoot[R]{\footnotesize\color{muted}\sffamily Kaicheng Dong \enspace | \enspace \thepage}
\titleformat{\section}{\large\sffamily\bfseries\color{accent}}{}{0pt}{}[\vspace{2pt}\color{rule}\titlerule]
\titlespacing*{\section}{0pt}{15pt}{8pt}
\newcommand{\rightdate}[1]{\hfill{\footnotesize\sffamily\color{muted} #1}}
\newcommand{\entrytitle}[2]{\needspace{3\baselineskip}\textbf{#1}\rightdate{#2}\par}
\newcommand{\subnote}[1]{{\color{muted}#1}\par}
\newcommand{\indented}[1]{\begin{list}{}{\setlength{\leftmargin}{1.05em}\setlength{\rightmargin}{0em}\setlength{\topsep}{0pt}\setlength{\itemsep}{0pt}\setlength{\parsep}{0pt}}\item #1\end{list}}
\setlist[itemize]{label=\textcolor{accent}{\textbullet},nosep,leftmargin=1.1em,topsep=1pt,parsep=0pt,itemsep=3pt}
\begin{document}
\begin{center}
{\fontsize{25}{29}\selectfont\sffamily\bfseries Kaicheng Dong}\par
\vspace{5pt}
{\sffamily\small
\href{mailto:12012612@mail.sustech.edu.cn}{12012612@mail.sustech.edu.cn}
\enspace\textcolor{rule}{|}\enspace
\href{https://dongkaicheng.github.io/}{dongkaicheng.github.io}
\enspace\textcolor{rule}{|}\enspace
\href{https://openmath-atlas.org/}{OpenMath}
}\par
\end{center}
\vspace{-7pt}
\section*{Research interests}
Arithmetic and Diophantine geometry, Arakelov geometry, adelic line bundles and heights, arithmetic intersection theory, and non-Archimedean geometry. Particular interests include positivity, curves and Jacobians, and Bogomolov-type problems.

\section*{Education}
\entrytitle{ALGANT Master's Programme in Mathematics (120 ECTS)}{2024--2026}
\subnote{University of Regensburg, Germany (2024--2025); University of Bordeaux, France (2025--2026)}
\indented{Completed the programme in 2026. Master's thesis supervised by Professor Pascal Autissier. Bordeaux M2: 60 ECTS, final result 15.85/20.}
\vspace{4pt}
\entrytitle{B.Sc. in Mathematics and Applied Mathematics}{2020--2024}
\subnote{Southern University of Science and Technology (SUSTech), Shenzhen, China}
\indented{Mathematics Honors Programme; GPA \textbf{3.72/4.00}; class rank \textbf{7/46}.}

\section*{Academic visit}
\entrytitle{Peking University, Beijing, China}{2026--2027}
\indented{Non-degree visiting student; host: Professor Xinyi Yuan.}

\section*{Preprint}
\entrytitle{\href{https://arxiv.org/abs/2611.09319}{\emph{One complex metric and geometric adelic nefness}}}{November 2026}
\subnote{\href{https://arxiv.org/abs/2611.09319}{arXiv:2611.09319}.}

\section*{Research experience}
\entrytitle{\href{https://github.com/DongKaicheng/thesis-supplementary-files/releases/download/v1.0/Dong_Kaicheng_Master_Thesis.pdf.pdf}{\emph{Admissible Pairing and Effective Bogomolov Results}}}{2026}
\subnote{Master's thesis, University of Bordeaux; supervisor: Professor Pascal Autissier. \href{https://github.com/DongKaicheng/thesis-supplementary-files/releases/download/v1.0/Dong_Kaicheng_Master_Thesis.pdf.pdf}{Thesis PDF}.}
\vspace{5pt}
\entrytitle{\emph{Serre's Conjecture II and Galois Cohomology}}{2023--2024}
\subnote{Undergraduate research project, SUSTech; supervisor: Professor Hu Yong.}


\section*{AI for Mathematics and research tools}
\entrytitle{Founder and maintainer of \href{https://openmath-atlas.org/}{OpenMath}}{2026--present}
\indented{Built an open mathematical research atlas for interconnected definitions, conjectures, proof dependencies, revision histories, and reader feedback; includes non-Archimedean and arithmetic geometry projects.}
\vspace{2pt}
\entrytitle{Creator of three \href{https://openmath-atlas.org/ai-and-math/}{downloadable mathematical research agents}}{2026}
\indented{\emph{Math Project Maintainer} (theorem records and dependency graphs); \emph{Math Research Publisher} (versioned research websites); \emph{Math Writing Editor} (mathematical manuscripts). The tools support source-linked logical diagrams and human review; they do not constitute independent proof certification.}

\clearpage
\section*{Talk}
\entrytitle{\emph{Arithmetic Bigness and Uniform Bogomolov-type Theory}}{2025}
ALGANT Alumni in China Conference (20 Years of ALGANT), University of Science and Technology of China.
\href{https://faculty.ustc.edu.cn/yqliang/files/ALGANTSlides/25DongKaicheng.pdf}{Slides}.

\section*{Selected mathematical coursework}
\entrytitle{University of Regensburg}{2024--2025}
\indented{Diophantine Geometry I (1.0), Diophantine Geometry II (1.3), Algebraic Number Theory I (1.3), Non-Archimedean \mbox{Banach} Algebras (1.0), Algebraic Topology I (1.0). \textit{German grading scale: 1.0 is the highest grade.}}
\vspace{3pt}
\entrytitle{University of Bordeaux}{2025--2026}
\indented{Algebraic Geometry (17.6/20), Advanced Analysis (16.5/20), Geometry (14.2/20), Advanced Geometry (14.0/20), Research Project (19.0/20).}
\vspace{3pt}
\entrytitle{Southern University of Science and Technology}{2020--2024}
\indented{Algebraic Geometry (95/100), Algebraic Curves (98/100), Complex Analysis (95/100), Real Analysis (94/100), Algebraic Topology (88/100), Elementary Number Theory (81/100).}

\section*{Seminars}
\entrytitle{Reading seminar on algebraic fibered surfaces}{2026}
Organizer. Topics included intersection theory on arithmetic surfaces, regular models, semistable reduction, and stable curves; based on Qing Liu's \emph{Algebraic Geometry and Arithmetic Curves}.
\vspace{4pt}
\entrytitle{Reading seminar on Berkovich differential forms and currents}{2026}
Participant and speaker. Seminar organized by Mingchen Xia; presented sections on local properties of Berkovich analytic spaces and morphisms.
\vspace{4pt}
\entrytitle{Further graduate seminar training}{2024--2025}
Regensburg: algebraic function fields, class field theory, Galois representations and $L$-functions, and \'etale cohomology.

\section*{Awards and schools}
\begin{itemize}
\item Erasmus+ Mobility Grant, European Union (2025--2026).
\item Mathematics Honors Programme, SUSTech (2020--2024).
\item Summer School on Arithmetic Geometry, Westlake University, China (2024).
\end{itemize}

\end{document}
